% 摘要专用页（官方第三条：含标题与关键词，不超过一页；电子版论文第一页即本页）
\thispagestyle{plain}
{\linespread{1.25}\selectfont
\begin{center}
  {\zihao{3}\bfseries 药材的烘干问题：圆柱非稳态传热传质\\耦合模型的建立、达标时间反演与收缩效应分析}
\end{center}

\vspace{0.4em}
\begin{center}\zihao{4}\bfseries 摘要\end{center}

\vspace{0.2em}

中药材干燥是产地加工的关键环节。本文建立一维轴对称圆柱非稳态传热传质耦合模型：控制方程取散度形式以兼容变物性，空间用元体平衡以严格守恒，时间推进分场选取；四问共用同一模型，仅物性、耦合方式、终止判据与几何假设不同。

针对问题一，题面给定常物性且两场解耦。温度方程用时间方向精确推进（矩阵指数与分段线性杜哈梅尔积分）使时间误差归零，含水率方程用后向欧拉全隐式配合追赶法，网格 0.25 mm、子步 0.1 s。1800 s 时中心温度 33.5753 ℃、表面 36.7855 ℃，含水率中心 2.5500 kg/kg 未变、表面降至 1.5104 kg/kg、降幅 40.8\%；步长 1 s 与 60 s 相差 3.1×10\textsuperscript{-11} ℃，印证误差归零。

针对问题二，物性升为含水率与温度的函数，两场强耦合。以交替推进求解，物性在每个内部子步按局部状态更新（子步 1/32 s），含水率侧配合欠松弛（0.7），界面变系数取含水率积分平均。3 h 时中心温度 49.8495 ℃、表面 49.9664 ℃，含水率中心 1.7662 kg/kg、表面 1.0081 kg/kg，较初值降约 30.7\% 与 60.5\%。题面未给汽化潜热，不含源项的代价为 0.07 至 0.10 ℃ 温降。

针对问题三，以全域含水率低于 0.15 kg/kg 为判据，事件驱动终止、子步级二分定位达标时刻，得烘干时长 57.53 h，带 ±0.08 h 离散误差带（两类误差带取大者）。中心含水率前 24 h 降约九成，其后 33.5 h 仅由 0.2378 降至 0.1500 kg/kg。低含水率端外推贡献占时长 43\% 至 66\%，先验为主观设定。

针对问题四，引入物料坐标把收缩域固定为计算域，守恒形式写能量方程。烘干时长 50.6523 h，离散误差带 ±0.0249 h，较固定半径加快约 61\%；净效应为物性变化延长 71.9 h 与收缩缩短 78.8 h，合计较问题三缩短 6.9 h。模型还揭示压缩温升，最高温度 54.63 ℃ 高于环境上限。

以十类检验支撑结论：误差带、解析对拍、独立实现互验、守恒残差 10\textsuperscript{-4} 至 10\textsuperscript{-2}、两族灵敏度与四组对照。

本文改进集中在数值格式与问题特化：时间方向精确推进取代加密步长逼近，使温度场时间误差归零；事件驱动终止配合子步级二分把达标时刻定位到 10\textsuperscript{-3} s；物料坐标与守恒形式配合，消除移动域插值误差与表面列收敛困难。本文为单一机理模型的四问递进求解，未做多模型融合。

\keywords{非稳态传热传质；圆柱坐标系有限体积法；后向欧拉隐式格式；参数反演；移动边界}
}
% ← 关闭「摘要内容行距 1.25」分组（正文仍 1.0 单倍）

% 摘要页边界：官方要求摘要单独一页（属**前置件边界**，非「章/节间人为断页」）
\clearpage
